\clearpage
\hypertarget{native-484-psyche}{}
\begin{center}
\rule{0.55\linewidth}{0.4pt}\\[2pt]
Astronomische Nachrichten. Band XXXV. Nr. 818. S. 17. 27. Juli 1852.\\[2pt]
\rule{0.55\linewidth}{0.4pt}
\end{center}
\begin{center}Elemente der \emph{Psyche} mit Bemerkungen zur Bestimmung der Bahn dieses Planeten.\end{center}

Weil die Bahn der Psyche so wenig geneigt ist, eignet sich eine Bestimmung aus drei Beobachtungen überhaupt nicht für genaue Ergebnisse. Das bestätigen auch die Elemente, die verschiedene Rechner auf diese Weise erhalten haben. Ich habe daher Herrn \textsc{Klinkerfues} veranlasst, eine neue Bestimmung auf vier Orte zu stützen und dazu die in der Th. M. C. C. entwickelte Methode anzuwenden, die bisher, wie es scheint, wenig oder gar nicht verwendet wurde. Ich lege seine Elemente hier bei, zusammen mit dem Vergleich mit allen Beobachtungen, die uns bekannt geworden sind.

Interessant ist die enge Übereinstimmung mit \textsc{Hind}s Schätzung vom 29. Januar. Damit ist nicht mehr daran zu zweifeln, dass der damals von \textsc{Hind} gesehene Stern derselbe ist wie der 7 Wochen später von \textsc{Gasparis} entdeckte Planet. Wenn Herr \textsc{Hind} Näheres zu seiner damaligen Aufzeichnung mitteilen wollte, insbesondere zur Stunde und dazu, für wie genau er seine Schätzung hält, könnte die Einbeziehung dieser Beobachtung in eine neue Berechnung die Genauigkeit des Ergebnisses wahrscheinlich noch merklich erhöhen.

\Needspace{17\baselineskip}
\begin{center}
Elemente der Psyche.\\
Epoche: März 31.0, mittlere Berliner Zeit.
\end{center}
\[
\text{Mittlere Länge}=151^\circ\ 32'\ 5''56
\]
\[
\left.\begin{array}{rcl}
\pi &=& 6\ 20\ 2.86 \\
\AscendingNode &=& 150\ 33\ 5.86
\end{array}\right\}\quad\text{Mittl. Aequin. 1852.0}
\]
\[
\begin{array}{rcl}
i &=& 3\ 2\ 16.06 \\
\varphi &=& 6\ 38\ 32.80 \\
\log e &=& 9.0632320 \\
\log a &=& 0.4693195 \\
\mu &=& 701''4993
\end{array}
\]

Die vorstehenden Elemente wurden aus 4 Normalorten vom 19. März, 14. April, 30. April und 14. Mai berechnet.

Zu den unten [hier nicht abgedruckten] Beobachtungen, die mit diesem System verglichen wurden, gehört auch die Schätzung von Herrn \textsc{Hind} vom 29. Januar. Da die Uhrzeit zu dieser Schätzung unbekannt ist, wurde dafür ungefähr Berliner Mitternacht angenommen. Zu dieser Zeit beträgt die stündliche Bewegung des Planeten $22''75$ nach Westen und $10''25$ nach Norden.

